% ============================================================
%  C: Como Programar -- 6ª Edição | Deitel & Deitel
%  Capítulo 9 -- Entrada/Saída Formatada em C
%  EXERCÍCIOS (9.4 -- 9.20)
%  Abordagem incremental: Caps. 1 a 9
% ============================================================
\documentclass[12pt,a4paper]{article}
\usepackage[utf8]{inputenc}
\usepackage[T1]{fontenc}
\usepackage[brazil]{babel}
\usepackage{geometry}
\geometry{top=2.5cm,bottom=2.5cm,left=3cm,right=3cm}
\usepackage{parskip}\setlength{\parskip}{6pt}\setlength{\parindent}{0pt}
\usepackage{xcolor,tcolorbox,enumitem,listings,fancyhdr,titlesec,hyperref,booktabs,array,amsmath}
\tcbuselibrary{skins,breakable}

\definecolor{deitelblue}{RGB}{0,70,127}
\definecolor{deitelgray}{RGB}{240,242,245}
\definecolor{deitelgreen}{RGB}{0,110,60}
\definecolor{deitellightblue}{RGB}{220,235,250}
\definecolor{deitelred}{RGB}{160,20,20}
\definecolor{deitellorange}{RGB}{200,90,0}

\newtcolorbox{boaprática}[1][]{colback=deitellightblue,colframe=deitelblue,
  fonttitle=\bfseries\small,title={Boa prática de programação},breakable,#1}
\newtcolorbox{erroco}[1][]{colback=red!5,colframe=deitelred,
  fonttitle=\bfseries\small,title={Erro comum de programação},breakable,#1}
\newtcolorbox{respostabox}[1][]{colback=deitelgray,colframe=deitelblue!60,
  fonttitle=\bfseries\small\color{deitelblue},title={Resposta detalhada},breakable,#1}
\newtcolorbox{enunciadoBox}[1][]{colback=white,colframe=deitelblue,
  fonttitle=\bfseries,title={#1},breakable}
\newtcolorbox{saida}[1][]{colback=black!88,colframe=deitelblue,
  fontupper=\ttfamily\small\color{green!70},
  title={\small\color{white}Saída do programa},breakable,#1}

\setlist[enumerate,1]{label=\textbf{\alph*)},leftmargin=2em}
\setlist[itemize]{leftmargin=2em}

\lstset{
  basicstyle=\ttfamily\small,backgroundcolor=\color{deitelgray},
  frame=single,framerule=0.4pt,rulecolor=\color{deitelblue!40},
  breaklines=true,showstringspaces=false,
  keywordstyle=\color{deitelblue}\bfseries,
  commentstyle=\color{deitelgreen}\itshape,
  stringstyle=\color{deitelred},
  language=C,numbers=left,numberstyle=\tiny\color{gray},
  xleftmargin=18pt,xrightmargin=5pt,stepnumber=1,
}

\pagestyle{fancy}\fancyhf{}
\fancyhead[L]{\small\color{deitelblue}\textbf{C: Como Programar -- 6ª ed. | Deitel \& Deitel}}
\fancyhead[R]{\small\color{deitelblue}Cap. 9 -- Exercícios}
\fancyfoot[C]{\small\thepage}
\renewcommand{\headrulewidth}{0.4pt}

\titleformat{\section}[block]{\large\bfseries\color{deitelblue}}{\thesection.}{0.5em}{}[\titlerule]
\titleformat{\subsection}[block]{\normalsize\bfseries\color{deitelblue!80}}{}{}{}

\hypersetup{colorlinks=true,linkcolor=deitelblue}

\begin{document}

\begin{titlepage}
  \centering\vspace*{2cm}
  {\color{deitelblue}\rule{\linewidth}{2pt}}\\[0.4cm]
  {\huge\bfseries\color{deitelblue}C: Como Programar}\\[0.2cm]
  {\large\color{deitelblue}6ª Edição $\cdot$ Paul Deitel \& Harvey Deitel}\\[0.2cm]
  {\color{deitelblue}\rule{\linewidth}{2pt}}\\[1cm]
  {\LARGE\bfseries Capítulo 9}\\[0.3cm]
  {\Large\color{deitelblue!80}Entrada/Saída Formatada em C}\\[0.8cm]
  \begin{tcolorbox}[colback=deitellightblue,colframe=deitelblue,width=0.85\linewidth]
    \centering{\large\bfseries Exercícios (9.4--9.20)}\\[0.2cm]
    {\normalsize Resolvidos passo a passo, com código completo}
  \end{tcolorbox}
\end{titlepage}

\tableofcontents\newpage

% ============================================================
\section{Exercício 9.4 -- Instruções \texttt{printf} e \texttt{scanf}}
% ============================================================

\begin{respostabox}
\begin{enumerate}[label=\textbf{\alph*)}]

\item \textbf{Inteiro sem sinal 40000 alinhado à esquerda em campo de 15:}
\begin{lstlisting}
printf( "%-15u\n", 40000 );
\end{lstlisting}

\item \textbf{Ler hexadecimal em \texttt{hex}:}
\begin{lstlisting}
scanf( "%x", &hex );
\end{lstlisting}

\item \textbf{Imprimir 200 com e sem sinal:}
\begin{lstlisting}
printf( "%d %u\n", 200, 200 );
\end{lstlisting}

\item \textbf{Imprimir 100 em hex precedido por 0x:}
\begin{lstlisting}
printf( "%#x\n", 100 );  /* imprime 0x64 */
\end{lstlisting}
A flag \texttt{\#} com \texttt{x} adiciona o prefixo \texttt{0x}.

\item \textbf{Ler caracteres em \texttt{s} até encontrar 'p':}
\begin{lstlisting}
scanf( "%[^p]", s );
\end{lstlisting}
O \texttt{\^{}} dentro do conjunto significa ``qualquer caractere \textit{exceto}''
os listados. Aqui: lê tudo até encontrar \texttt{p}.

\item \textbf{Imprimir 1.234 em campo de 9 com zeros iniciais:}
\begin{lstlisting}
printf( "%09.3f\n", 1.234 );  /* imprime 00001.234 */
\end{lstlisting}
A flag \texttt{0} preenche com zeros à esquerda; \texttt{.3} mantém 3 decimais.

\item \textbf{Ler \texttt{hh:mm:ss} com supressão dos \texttt{:}:}
\begin{lstlisting}
scanf( "%d%*c%d%*c%d", &hora, &minuto, &segundo );
\end{lstlisting}
\texttt{\%*c} lê um caractere e descarta-o (supressão de atribuição).

\item \textbf{Ler ``caracteres'' (com aspas) e armazenar em \texttt{s}:}
\begin{lstlisting}
scanf( "\"%[^\"]\"", s );
\end{lstlisting}
Lê literalmente as aspas iniciais e finais, e captura em \texttt{s} todos os
caracteres entre elas.

\item \textbf{Ler \texttt{hh:mm:ss} sem supressão:}
\begin{lstlisting}
scanf( "%d:%d:%d", &hora, &minuto, &segundo );
\end{lstlisting}
Os \texttt{:} na string de formato \textit{esperam} encontrar \texttt{:} na entrada
(que são descartados).

\end{enumerate}
\end{respostabox}

\newpage

% ============================================================
\section{Exercício 9.5 -- Saídas das Instruções}
% ============================================================

\begin{respostabox}
\begin{enumerate}[label=\textbf{\alph*)}]

\item \textbf{\texttt{printf( "\%-10d\textbackslash n", 10000 );}}\\
\textbf{Saída:} \texttt{10000\textvisiblespace\textvisiblespace\textvisiblespace\textvisiblespace\textvisiblespace}\\
(10000 ocupa 5 caracteres, com 5 espaços à direita até totalizar 10.)

\item \textbf{\texttt{printf( "\%c\textbackslash n", "Isto e uma string" );}}\\
\textbf{ERRO:} \texttt{\%c} espera um \texttt{char}, mas recebeu um ponteiro para
\texttt{char}. Comportamento indefinido -- provavelmente imprime um caractere
``estranho'' ou causa falha.\\
\textbf{Correção:} \texttt{printf( "\%s\textbackslash n", "Isto e uma string" );}

\item \textbf{\texttt{printf( "\%*.*lf\textbackslash n", 8, 3, 1024.987654 );}}\\
Largura dinâmica = 8, precisão dinâmica = 3.\\
\textbf{Saída:} \texttt{1024.988} \quad (8 caracteres, 3 decimais arredondados)

\item \textbf{\texttt{printf( "\%\#o\textbackslash n\%\#X\textbackslash n\%\#e\textbackslash n", 17, 17, 1008.83689 );}}\\
\textbf{Saída:}
\begin{verbatim}
021
0X11
1.008837e+03
\end{verbatim}

\item \textbf{\texttt{printf( "\% ld\textbackslash n\%+ld\textbackslash n", 1000000, 1000000 );}}\\
Flag espaço: insere espaço antes de positivos. Flag \texttt{+}: insere \texttt{+}.\\
\textbf{Saída:}
\begin{verbatim}
 1000000
+1000000
\end{verbatim}

\item \textbf{\texttt{printf( "\%10.2E\textbackslash n", 444.93738 );}}\\
\textbf{Saída:} \texttt{\textvisiblespace\textvisiblespace 4.45E+02} \quad (largura 10, 2 decimais, E maiúsculo)

\item \textbf{\texttt{printf( "\%10.2g\textbackslash n", 444.93738 );}}\\
\texttt{\%g} usa o formato mais compacto. Com precisão 2 (2 dígitos significativos):\\
\textbf{Saída:} \texttt{\textvisiblespace\textvisiblespace\textvisiblespace\textvisiblespace\textvisiblespace 4.4e+02}

\item \textbf{\texttt{printf( "\%d\textbackslash n", 10.987 );}}\\
\textbf{ERRO:} \texttt{\%d} espera \texttt{int}, mas recebeu \texttt{double}.
Comportamento indefinido (em sistemas modernos pode imprimir lixo ou crashar).\\
\textbf{Correção:}
\begin{lstlisting}
printf( "%d\n", ( int )10.987 );    /* cast explicito */
printf( "%f\n", 10.987 );           /* ou use %f      */
\end{lstlisting}
\end{enumerate}
\end{respostabox}

\newpage

% ============================================================
\section{Exercício 9.6 -- Erros nos Segmentos}
% ============================================================

\begin{respostabox}
\begin{enumerate}[label=\textbf{\alph*)}]

\item \textbf{\texttt{printf( "\%s\textbackslash n", 'Feliz aniversario' );}}\\
\textbf{Erro:} String entre aspas simples. Strings exigem aspas duplas.\\
\textbf{Correção:} \texttt{printf( "\%s\textbackslash n", "Feliz aniversario" );}

\item \textbf{\texttt{printf( "\%c\textbackslash n", 'Ola' );}}\\
\textbf{Erro:} Aspas simples com múltiplos caracteres (\texttt{'Ola'}) gera apenas
um \texttt{int} com valor implementação-definido.\\
\textbf{Correção:} para imprimir apenas um caractere, use \texttt{'O'} (um único
caractere); para a palavra, use \texttt{\%s} com \texttt{"Ola"}.

\item \textbf{\texttt{printf( "\%c\textbackslash n", "Isso e uma string" );}}\\
\textbf{Erro:} \texttt{\%c} espera caractere, recebeu string.\\
\textbf{Correção:} use \texttt{\%s}, ou para imprimir apenas o primeiro caractere:
\begin{lstlisting}
printf( "%c\n", "Isso e uma string"[ 0 ] );
\end{lstlisting}

\item \textbf{\texttt{printf( ""\%s"", "Boa viagem" );}}\\
\textbf{Erro:} Aspas duplicadas dentro da string de formato sem escape.\\
\textbf{Correção:}
\begin{lstlisting}
printf( "\"%s\"", "Boa viagem" );  /* imprime "Boa viagem" */
\end{lstlisting}

\item \textbf{\texttt{char day[] = "Domingo"; printf( "\%s\textbackslash n", day[ 3 ] );}}\\
\textbf{Erro:} \texttt{day[3]} é um \texttt{char} (no caso, \texttt{'i'}), não um
ponteiro. \texttt{\%s} espera \texttt{char *}.\\
\textbf{Correção:}
\begin{lstlisting}
printf( "%s\n", &day[ 3 ] );  /* imprime "ingo" a partir do 4o char */
\end{lstlisting}

\item \textbf{\texttt{printf( 'Digite seu nome: ' );}}\\
\textbf{Erro:} String entre aspas simples.\\
\textbf{Correção:} \texttt{printf( "Digite seu nome: " );}

\item \textbf{\texttt{printf( \%f, 123.456 );}}\\
\textbf{Erro:} String de formato sem aspas.\\
\textbf{Correção:} \texttt{printf( "\%f", 123.456 );}

\item \textbf{\texttt{printf( "\%s\%s\textbackslash n", 'O', 'K' );}}\\
\textbf{Erro:} \texttt{\%s} espera strings, recebeu caracteres em aspas simples.\\
\textbf{Correção:}
\begin{lstlisting}
printf( "%c%c\n", 'O', 'K' );    /* dois caracteres */
printf( "%s%s\n", "O", "K" );    /* duas strings    */
\end{lstlisting}

\item \textbf{\texttt{char s[10]; scanf( "\%c", s[ 7 ] );}}\\
\textbf{Erro:} Falta o operador de endereço \texttt{\&}. \texttt{scanf} precisa de
ponteiro para onde armazenar.\\
\textbf{Correção:} \texttt{scanf( "\%c", \&s[ 7 ] );}

\end{enumerate}
\end{respostabox}

\newpage

% ============================================================
\section{Exercício 9.7 -- Diferenças entre \texttt{\%d} e \texttt{\%i}}
% ============================================================

\begin{respostabox}
\begin{lstlisting}
/* Ex9.7: d_vs_i.c */
#include <stdio.h>

int main( void )
{
    int x, y;

    printf( "Digite dois inteiros: " );
    scanf( "%i%d", &x, &y );
    printf( "Lidos: x = %d, y = %d\n", x, y );

    return 0;
}
\end{lstlisting}

\textbf{Diferença chave:} \texttt{\%d} sempre lê em \textit{decimal}. Já \texttt{\%i}
\textit{detecta} a base automaticamente baseando-se no prefixo:
\begin{itemize}[noitemsep]
  \item \texttt{0x} ou \texttt{0X} $\to$ hexadecimal
  \item \texttt{0} (zero à esquerda) $\to$ octal
  \item Sem prefixo $\to$ decimal
\end{itemize}

\textbf{Testes com os conjuntos sugeridos:}
\begin{center}
\renewcommand{\arraystretch}{1.3}
\begin{tabular}{l|l|l}
\toprule
\textbf{Entrada} & \texttt{\%i} (x) & \texttt{\%d} (y) \\ \midrule
\texttt{10 10}     & 10  & 10  \\
\texttt{-10 -10}   & -10 & -10 \\
\texttt{010 010}   & 8 (octal) & 10 (decimal, ignora 0) \\
\texttt{0x10 0x10} & 16 (hex)  & 0 (lê 0, para no 'x') \\
\bottomrule
\end{tabular}
\end{center}

\textbf{Conclusão:} \texttt{\%i} é mais ``inteligente''. Útil quando o usuário pode
inserir números em bases diferentes. Para uso geral, \texttt{\%d} é mais previsível.
\end{respostabox}

% ============================================================
\section{Exercício 9.8 -- \texttt{\%n} para Contar Caracteres}
% ============================================================

\begin{respostabox}
\begin{lstlisting}
/* Ex9.8: contador_chars.c */
#include <stdio.h>
#include <stdlib.h>
#include <time.h>
#define N 10

int main( void )
{
    int number[ N ];
    int impressos;       /* caracteres impressos por valor */
    int total = 0;       /* acumulado */
    int i;

    srand( time( NULL ) );

    /* Inicializa array */
    for ( i = 0; i < N; ++i ) {
        number[ i ] = 1 + rand() % 1000;
    }

    printf( "%-10s%-15s%s\n", "Valor", "Chars", "Total acumulado" );

    for ( i = 0; i < N; ++i ) {
        printf( "%-10d%n", number[ i ], &impressos );

        /* impressos contem o numero de chars escritos ATE %n */
        total += impressos - 0;  /* subtrai o tamanho ja somado */

        /* Re-imprime o valor + total ate aqui */
        total = 10;  /* largura do campo */
        printf( "Total acumulado: %d\n", total * ( i + 1 ) );
    }

    return 0;
}
\end{lstlisting}

\textit{Nota:} \texttt{\%n} é um especificador especial que \textit{não imprime
nada}, mas armazena no inteiro apontado pelo argumento o número de caracteres
já impressos até aquele ponto na string de saída.

\begin{erroco}
  \texttt{\%n} é considerado \textbf{perigoso} em código moderno. Sistemas com
  \textit{format string vulnerability} podem ser explorados via \texttt{\%n}.
  Por isso, muitos compiladores e bibliotecas (incluindo Windows) o desabilitam
  por padrão. Use com cuidado, apenas em contextos educacionais.
\end{erroco}
\end{respostabox}

\newpage

% ============================================================
\section{Exercício 9.9 -- Ponteiros como Inteiros}
% ============================================================

\begin{respostabox}
\begin{lstlisting}
/* Ex9.9: ponteiros_int.c */
#include <stdio.h>

int main( void )
{
    int x = 42;
    int *ptr = &x;

    printf( "%%p:  %p\n", ( void * )ptr );

    /* Possivelmente OK (depende do compilador) */
    printf( "%%d:  %d\n", ( int )( long )ptr );
    printf( "%%u:  %u\n", ( unsigned )( long )ptr );

    /* Em sistemas 64-bit, int causa truncamento */
    printf( "%%lu: %lu\n", ( unsigned long )ptr );
    printf( "%%lx: %lx\n", ( unsigned long )ptr );
    printf( "%%o:  %o\n", ( unsigned )( long )ptr );

    return 0;
}
\end{lstlisting}

\textbf{Resultados esperados:}
\begin{itemize}[noitemsep]
  \item \textbf{\texttt{\%p}}: formato definido pela implementação. Geralmente
  em hexadecimal, e.g.\ \texttt{0x7fff5fbff8ac} em sistemas Linux/macOS.
  \item \textbf{\texttt{\%d}, \texttt{\%u}}: em sistemas 64-bit, gera valor truncado
  (perda de bits altos) -- \textbf{valores estranhos}.
  \item \textbf{\texttt{\%lu}, \texttt{\%lx}}: representação completa em decimal/hex.
  \item \textbf{\texttt{\%o}}: representação octal -- valores ``estranhos'', muito longos.
\end{itemize}

\textbf{Erros possíveis:} sem cast adequado, o compilador moderno emite avisos.
Casts inadequados podem levar a comportamento indefinido em arquiteturas onde
\texttt{int} é menor que \texttt{void *}.
\end{respostabox}

% ============================================================
\section{Exercício 9.10 -- Larguras de Campo Diversas}
% ============================================================

\begin{respostabox}
\begin{lstlisting}
/* Ex9.10: larguras_diversas.c */
#include <stdio.h>

int main( void )
{
    int   inteiro = 12345;
    float real    = 1.2345f;
    int   largura;

    printf( "%-10s | %-10s\n", "Largura", "Saida" );
    printf( "-----------+--------------\n" );

    for ( largura = 1; largura <= 10; ++largura ) {
        printf( "%-10d | %*d\n", largura, largura, inteiro );
    }

    printf( "\n%-10s | %-15s\n", "Largura", "Saida float" );
    printf( "-----------+--------------------\n" );

    for ( largura = 1; largura <= 10; ++largura ) {
        printf( "%-10d | %*.4f\n", largura, largura, real );
    }

    return 0;
}
\end{lstlisting}

\textbf{O que acontece com largura insuficiente?} \textbf{C \textit{nunca} trunca o número}.
Quando o valor não cabe na largura especificada, ele é impresso por completo,
``estourando'' o campo. Ou seja, a largura é um \textit{mínimo}, não um máximo.

\textbf{Exemplo:} \texttt{printf("\%3d", 12345)} imprime \texttt{12345} (5 chars), não \texttt{123}.
\end{respostabox}

\newpage

% ============================================================
\section{Exercício 9.11 -- Arredondamento}
% ============================================================

\begin{respostabox}
\begin{lstlisting}
/* Ex9.11: arredondamento.c */
#include <stdio.h>

int main( void )
{
    double valor = 100.453627;

    printf( "Original:           %.6f\n", valor );
    printf( "Inteiro (%%.0f):     %.0f\n", valor );
    printf( "Decimo (%%.1f):      %.1f\n", valor );
    printf( "Centesimo (%%.2f):   %.2f\n", valor );
    printf( "Milesimo (%%.3f):    %.3f\n", valor );
    printf( "Decimo milesimo:    %.4f\n", valor );

    return 0;
}
\end{lstlisting}

\textbf{Saída:}
\begin{saida}
Original:           100.453627\\
Inteiro (\%.0f):     100\\
Decimo (\%.1f):      100.5\\
Centesimo (\%.2f):   100.45\\
Milesimo (\%.3f):    100.454\\
Decimo milesimo:    100.4536
\end{saida}

\textbf{Note:} \texttt{printf} arredonda matematicamente, não trunca.
\texttt{100.4536} arredondado para a quarta casa: o quinto dígito é 2, então mantém.
Mas \texttt{100.453627} arredondado para 3 casas: o quarto dígito é 6 ($\geq 5$),
arredonda para cima $\to$ \texttt{100.454}.
\end{respostabox}

% ============================================================
\section{Exercício 9.12 -- String em Campo do Dobro}
% ============================================================

\begin{respostabox}
\begin{lstlisting}
/* Ex9.12: string_dobro.c */
#include <stdio.h>
#include <string.h>

int main( void )
{
    char texto[ 100 ];
    int  tam, larguraDobro;

    printf( "Digite uma string: " );
    fgets( texto, 100, stdin );
    texto[ strcspn( texto, "\n" ) ] = '\0';

    tam = strlen( texto );
    larguraDobro = tam * 2;

    printf( "Comprimento: %d\n", tam );
    printf( "Dobrado (%d chars): [%*s]\n",
            larguraDobro, larguraDobro, texto );

    return 0;
}
\end{lstlisting}

\textbf{Exemplo:} se a entrada for ``Ola'' (3 chars), o campo será de 6 caracteres
e a saída terá 3 espaços à esquerda + ``Ola''.
\end{respostabox}

\newpage

% ============================================================
\section{Exercício 9.13 -- Tabela Fahrenheit/Celsius}
% ============================================================

\begin{respostabox}
\begin{lstlisting}
/* Ex9.13: tabela_temperaturas.c */
#include <stdio.h>

int main( void )
{
    int fahrenheit;
    double celsius;

    printf( "%10s%10s\n", "Fahr", "Celsius" );
    printf( "%10s%10s\n", "----", "-------" );

    for ( fahrenheit = 0; fahrenheit <= 212; ++fahrenheit ) {
        celsius = 5.0 / 9.0 * ( fahrenheit - 32 );
        printf( "%10d%+10.3f\n", fahrenheit, celsius );
    }

    return 0;
}
\end{lstlisting}

\textbf{Saída parcial:}
\begin{saida}
\textvisiblespace\textvisiblespace\textvisiblespace\textvisiblespace\textvisiblespace\textvisiblespace Fahr\textvisiblespace\textvisiblespace\textvisiblespace Celsius\\
\textvisiblespace\textvisiblespace\textvisiblespace\textvisiblespace\textvisiblespace\textvisiblespace ----\textvisiblespace\textvisiblespace\textvisiblespace -------\\
\textvisiblespace\textvisiblespace\textvisiblespace\textvisiblespace\textvisiblespace\textvisiblespace\textvisiblespace\textvisiblespace\textvisiblespace 0\textvisiblespace\textvisiblespace -17.778\\
\textvisiblespace\textvisiblespace\textvisiblespace\textvisiblespace\textvisiblespace\textvisiblespace\textvisiblespace\textvisiblespace\textvisiblespace 1\textvisiblespace\textvisiblespace -17.222\\
\ldots\\
\textvisiblespace\textvisiblespace\textvisiblespace\textvisiblespace\textvisiblespace\textvisiblespace\textvisiblespace\textvisiblespace 32\textvisiblespace\textvisiblespace\textvisiblespace +0.000\\
\textvisiblespace\textvisiblespace\textvisiblespace\textvisiblespace\textvisiblespace\textvisiblespace\textvisiblespace 100\textvisiblespace\textvisiblespace +37.778\\
\textvisiblespace\textvisiblespace\textvisiblespace\textvisiblespace\textvisiblespace\textvisiblespace\textvisiblespace 212\textvisiblespace +100.000
\end{saida}

\textbf{Flags-chave:} \texttt{\%+10.3f} -- sinal sempre exibido (\texttt{+}),
largura 10 alinhada à direita, 3 dígitos decimais.
\end{respostabox}

% ============================================================
\section{Exercício 9.14 -- Sequências de Escape}
% ============================================================

\begin{respostabox}
\begin{lstlisting}
/* Ex9.14: escapes.c */
#include <stdio.h>

int main( void )
{
    printf( "\\n: " ); printf( "[A\nB]\n" );
    printf( "\\t: [A\tB]\n" );
    printf( "\\a (alerta sonoro): [A\aB]\n" );
    printf( "\\b (backspace):    [AX\bB]\n" );
    printf( "\\r (retorno):      [ABC\rXY]\n" );
    printf( "\\f (form feed):    [A\fB]\n" );
    printf( "\\v (tab vertical): [A\vB]\n" );
    printf( "\\\\: [\\]\n" );
    printf( "\\\": [\"]\n" );
    printf( "\\?: [\?]\n" );
    printf( "\\0: [A\0B] (string termina no \\0)\n" );

    return 0;
}
\end{lstlisting}

\textbf{Comportamentos:}
\begin{itemize}[noitemsep]
  \item \texttt{\textbackslash n}: nova linha (cursor para próxima linha)
  \item \texttt{\textbackslash t}: tab horizontal (próxima posição de tab)
  \item \texttt{\textbackslash a}: alerta (geralmente beep)
  \item \texttt{\textbackslash b}: backspace (cursor recua 1 posição)
  \item \texttt{\textbackslash r}: carriage return (cursor volta ao início da linha)
  \item \texttt{\textbackslash f}: form feed (avança para próxima página em impressoras)
  \item \texttt{\textbackslash v}: tab vertical
  \item \texttt{\textbackslash \textbackslash}, \texttt{\textbackslash "}, \texttt{\textbackslash ?}: caracteres literais escapados
\end{itemize}
\end{respostabox}

\newpage

% ============================================================
\section{Exercício 9.15 -- ? Literal sem Escape}
% ============================================================

\begin{respostabox}
\begin{lstlisting}
/* Ex9.15: interrogacao.c */
#include <stdio.h>

int main( void )
{
    printf( "Com escape: \?\n" );
    printf( "Sem escape: ?\n" );
    return 0;
}
\end{lstlisting}

\textbf{Resposta:} \textbf{Sim}, o caractere \texttt{?} pode ser impresso diretamente
como caractere literal em uma string \texttt{printf}. As duas formas (\texttt{?}
e \texttt{\textbackslash ?}) produzem o mesmo resultado.

\textbf{Por que existe \texttt{\textbackslash ?} então?} Existe para evitar
\textit{trigraphs} (sequências de 3 caracteres que começam com \texttt{??}, usadas
em terminais antigos para representar caracteres especiais). Por exemplo,
\texttt{??/} era equivalente a \texttt{\textbackslash}. Em código moderno, trigraphs
foram removidos do padrão C23, mas \texttt{\textbackslash ?} ainda funciona.
\end{respostabox}

% ============================================================
\section{Exercício 9.16 -- Ler 437 com Todos os Especificadores}
% ============================================================

\begin{respostabox}
\begin{lstlisting}
/* Ex9.16: ler_437.c */
#include <stdio.h>

int main( void )
{
    int valor;

    printf( "Digite 437 (decimal): " );
    scanf( "%d", &valor );
    printf( "  %%d=%d  %%i=%d  %%o=%o  %%u=%u  %%x=%x\n",
            valor, valor, valor, valor, valor );

    printf( "Digite 437 (octal, ex: 665): " );
    scanf( "%o", &valor );
    printf( "  %%d=%d  %%o=%o\n", valor, valor );

    printf( "Digite 437 (hex, ex: 1b5): " );
    scanf( "%x", &valor );
    printf( "  %%d=%d  %%x=%x\n", valor, valor );

    return 0;
}
\end{lstlisting}

\textbf{Conversões equivalentes a 437:}
\begin{itemize}[noitemsep]
  \item Decimal: 437
  \item Octal: 665 (porque $6 \times 64 + 6 \times 8 + 5 = 437$)
  \item Hex: 0x1B5 (porque $1 \times 256 + 11 \times 16 + 5 = 437$)
\end{itemize}
\end{respostabox}

% ============================================================
\section{Exercício 9.17 -- Ler 1.2345 com \texttt{e}, \texttt{f}, \texttt{g}}
% ============================================================

\begin{respostabox}
\begin{lstlisting}
/* Ex9.17: ler_floats.c */
#include <stdio.h>

int main( void )
{
    float a, b, c;

    printf( "Digite 1.2345 em formato decimal: " );
    scanf( "%f", &a );

    printf( "Digite 1.2345 em formato exponencial (ex: 1.2345e0): " );
    scanf( "%e", &b );

    printf( "Digite 1.2345 em formato geral: " );
    scanf( "%g", &c );

    printf( "\nValores lidos: a=%f, b=%f, c=%f\n", a, b, c );

    return 0;
}
\end{lstlisting}

\textbf{Conclusão:} \texttt{\%e}, \texttt{\%f} e \texttt{\%g} no \texttt{scanf} são
\textit{equivalentes para leitura} -- todos aceitam tanto notação decimal quanto
científica. A diferença está apenas na \texttt{printf} (formato de saída).
\end{respostabox}

\newpage

% ============================================================
\section{Exercício 9.18 -- Strings com Aspas}
% ============================================================

\begin{respostabox}
\begin{lstlisting}
/* Ex9.18: aspas.c */
#include <stdio.h>

int main( void )
{
    char a[ 50 ], b[ 50 ], c[ 50 ];

    printf( "Digite 3 entradas: suzy \"suzy\" 'suzy'\n" );
    scanf( "%s%s%s", a, b, c );

    printf( "a = [%s]\n", a );
    printf( "b = [%s]\n", b );
    printf( "c = [%s]\n", c );

    return 0;
}
\end{lstlisting}

\textbf{Resultado:} C \textbf{não} ignora aspas -- elas são lidas como parte da
string. \texttt{\%s} para no primeiro espaço em branco. Então:
\begin{itemize}[noitemsep]
  \item \texttt{a = "suzy"} (sem aspas)
  \item \texttt{b = "\textbackslash"suzy\textbackslash""} (com as aspas duplas)
  \item \texttt{c = "'suzy'"} (com as aspas simples)
\end{itemize}

\textbf{Como ignorar aspas?} Usar \texttt{scanf("\textbackslash"\%[\textasciicircum\textbackslash"]\textbackslash"", b);}
que lê literalmente as aspas e captura o conteúdo entre elas.
\end{respostabox}

% ============================================================
\section{Exercício 9.19 -- \texttt{'?'} como Caractere}
% ============================================================

\begin{respostabox}
\begin{lstlisting}
/* Ex9.19: caractere_interrogacao.c */
#include <stdio.h>

int main( void )
{
    printf( "Constante de caractere '?': [%c]\n", '?' );
    printf( "Sequencia de escape  '\\?':  [%c]\n", '\?' );
    return 0;
}
\end{lstlisting}

\textbf{Resultado:} Ambas as formas (\texttt{'?'} e \texttt{'\textbackslash ?'})
funcionam perfeitamente e produzem o mesmo caractere de interrogação na saída.
A diferença é apenas estilística; \texttt{'\textbackslash ?'} historicamente protegia
contra trigraphs (hoje obsoleto).
\end{respostabox}

\newpage

% ============================================================
\section{Exercício 9.20 -- \texttt{\%g} com Precisões 1 a 9}
% ============================================================

\begin{respostabox}
\begin{lstlisting}
/* Ex9.20: g_precisoes.c */
#include <stdio.h>

int main( void )
{
    double valor = 9876.12345;
    int  prec;

    printf( "Valor: %f\n\n", valor );
    printf( "%-10s%s\n", "Precisao", "Saida com %g" );
    printf( "---------------------\n" );

    for ( prec = 1; prec <= 9; ++prec ) {
        printf( "%-10d%.*g\n", prec, prec, valor );
    }

    return 0;
}
\end{lstlisting}

\textbf{Saída esperada:}
\begin{saida}
Valor: 9876.123450\\
\\
Precisao  Saida com \%g\\
---------------------\\
1\textvisiblespace\textvisiblespace\textvisiblespace\textvisiblespace\textvisiblespace\textvisiblespace\textvisiblespace\textvisiblespace\textvisiblespace 1e+04\\
2\textvisiblespace\textvisiblespace\textvisiblespace\textvisiblespace\textvisiblespace\textvisiblespace\textvisiblespace\textvisiblespace\textvisiblespace 9.9e+03\\
3\textvisiblespace\textvisiblespace\textvisiblespace\textvisiblespace\textvisiblespace\textvisiblespace\textvisiblespace\textvisiblespace\textvisiblespace 9.88e+03\\
4\textvisiblespace\textvisiblespace\textvisiblespace\textvisiblespace\textvisiblespace\textvisiblespace\textvisiblespace\textvisiblespace\textvisiblespace 9876\\
5\textvisiblespace\textvisiblespace\textvisiblespace\textvisiblespace\textvisiblespace\textvisiblespace\textvisiblespace\textvisiblespace\textvisiblespace 9876.1\\
6\textvisiblespace\textvisiblespace\textvisiblespace\textvisiblespace\textvisiblespace\textvisiblespace\textvisiblespace\textvisiblespace\textvisiblespace 9876.12\\
7\textvisiblespace\textvisiblespace\textvisiblespace\textvisiblespace\textvisiblespace\textvisiblespace\textvisiblespace\textvisiblespace\textvisiblespace 9876.123\\
8\textvisiblespace\textvisiblespace\textvisiblespace\textvisiblespace\textvisiblespace\textvisiblespace\textvisiblespace\textvisiblespace\textvisiblespace 9876.1234\\
9\textvisiblespace\textvisiblespace\textvisiblespace\textvisiblespace\textvisiblespace\textvisiblespace\textvisiblespace\textvisiblespace\textvisiblespace 9876.12345
\end{saida}

\textbf{Comportamento do \texttt{\%g}:} a ``precisão'' em \texttt{\%g} significa
\textbf{dígitos significativos} (não casas decimais). O formato escolhe
automaticamente entre \texttt{\%e} e \texttt{\%f}, escolhendo o mais compacto:
\begin{itemize}[noitemsep]
  \item Precisão $\leq 4$: usa notação científica (\texttt{\%e}).
  \item Precisão $> 4$: usa notação decimal padrão (\texttt{\%f}).
  \item Zeros à direita são suprimidos (diferente de \texttt{\%f}).
\end{itemize}

\begin{boaprática}
  \texttt{\%g} é excelente quando você não sabe a magnitude dos valores que
  serão impressos. Mostra valores pequenos em decimal e valores muito grandes
  ou muito pequenos em notação científica -- a melhor representação compacta
  para cada caso.
\end{boaprática}
\end{respostabox}

\newpage

% ============================================================
\section*{Reflexão Final -- Capítulo 9}

\begin{tcolorbox}[colback=deitellightblue,colframe=deitelblue,
  title={\bfseries A arte da E/S formatada}]

O Capítulo~9 aprofunda algo que aparentemente já conhecíamos do Capítulo~2:
\texttt{printf} e \texttt{scanf}. Mas o que parecia simples revelou-se uma
\textbf{mini-linguagem} dentro de C, com dezenas de especificadores, flags,
modificadores e nuances.

\textbf{Domínio dos pontos-chave:}
\begin{itemize}[noitemsep]
  \item Saber qual especificador usar para cada tipo (\texttt{\%d}, \texttt{\%lf}, \texttt{\%s}, \texttt{\%p}, \ldots).
  \item Controlar largura, precisão e alinhamento para criar tabelas legíveis.
  \item Usar conjunto de varredura e supressão para parsear entradas complexas.
  \item Combinar flags (\texttt{+}, \texttt{-}, \texttt{0}, \texttt{\#}) com elegância.
  \item Reconhecer armadilhas: \texttt{\%d} com \texttt{double}, falta de \texttt{\&}
  no \texttt{scanf}, \texttt{\%c} com strings.
\end{itemize}

\textbf{Por que isso importa em projetos reais?} Logs, relatórios financeiros,
arquivos CSV, dashboards, parsers de configuração -- todos dependem de E/S formatada
com precisão. Programas amadores são reconhecíveis pela formatação descuidada;
programas profissionais primam por saídas limpas, alinhadas e informativas.

\textbf{Próximos passos:} no Cap.~10 chegaremos a \texttt{struct} -- agrupando
dados heterogêneos; no Cap.~11, arquivos -- aplicando tudo o que aprendemos
até agora a entrada/saída persistente em disco.

\end{tcolorbox}

\end{document}
